\documentclass{article}

% Recommended, but optional, packages for figures and better typesetting:
\usepackage[nopatch=footnote]{microtype}
\usepackage{graphicx}
\usepackage{subcaption}
\usepackage{booktabs} % for professional tables

% hyperref makes hyperlinks in the resulting PDF.
% If your build breaks (sometimes temporarily if a hyperlink spans a page)
% please comment out the following usepackage line and replace
% \usepackage{icml2026} with \usepackage[nohyperref]{icml2026} above.
\usepackage{hyperref}


% Attempt to make hyperref and algorithmic work together better:
\newcommand{\theHalgorithm}{\arabic{algorithm}}

% Use the following line for the initial blind version submitted for review:
\usepackage{icml2026}

% For preprint, use
% \usepackage[preprint]{icml2026}

% If accepted, instead use the following line for the camera-ready submission:
% \usepackage[accepted]{icml2026}

\usepackage{amsmath}
\usepackage{amssymb}
\usepackage{mathtools}
\usepackage{amsthm}


% if you use cleveref..
\usepackage[capitalize,noabbrev]{cleveref}

%%%%%%%%%%%%%%%%%%%%%%%%%%%%%%%%
% THEOREMS
%%%%%%%%%%%%%%%%%%%%%%%%%%%%%%%%
\theoremstyle{plain}
\newtheorem{theorem}{Theorem}[section]
\newtheorem{proposition}[theorem]{Proposition}
\newtheorem{lemma}[theorem]{Lemma}
\newtheorem{corollary}[theorem]{Corollary}
\theoremstyle{definition}
\newtheorem{definition}[theorem]{Definition}
\newtheorem{assumption}[theorem]{Assumption}
\theoremstyle{remark}
\newtheorem{remark}[theorem]{Remark}

% Todonotes is useful during development; simply uncomment the next line
%    and comment out the line below the next line to turn off comments
\usepackage[disable,textsize=tiny]{todonotes}
%\usepackage[textsize=tiny]{todonotes}

% The \icmltitle you define below is probably too long as a header.
% Therefore, a short form for the running title is supplied here:
\icmltitlerunning{Position: Don't Imprint Flawed Human Preferences Into AI}

\begin{document}

\twocolumn[
  \icmltitle{Position: Don't Imprint Flawed Human Preferences Into AI}

  % It is OKAY to include author information, even for blind submissions: the
  % style file will automatically remove it for you unless you've provided
  % the [accepted] option to the icml2026 package.
  \icmlsetsymbol{equal}{*}

  \begin{icmlauthorlist}
    \icmlauthor{Firstname1 Lastname1}{equal,yyy}
    \icmlauthor{Firstname2 Lastname2}{equal,comp}
  \end{icmlauthorlist}

  \icmlaffiliation{yyy}{Department of XXX, University of YYY, Location, Country}
  \icmlaffiliation{comp}{Company Name, Location, Country}

  \icmlcorrespondingauthor{Firstname1 Lastname1}{first1.last1@xxx.edu}
  \icmlcorrespondingauthor{Firstname2 Lastname2}{first2.last2@www.uk}

  \icmlkeywords{AI alignment, pluralistic alignment, RLHF, machine ethics, position paper}

  \vskip 0.3in
]

% this must go after the closing bracket ] following \twocolumn[ ...
\printAffiliationsAndNotice{}  % no special notice (required even if empty)

\begin{abstract}
  This position paper argues that aligning AI to aggregated human preferences is the wrong target. With current technology and modest expense, we can train AIs to share the values of a Silicon Valley techno-optimist, a degrowth environmentalist, a national-conservative culture warrior, a single-party state cadre, or a devout religious traditionalist. \emph{We should not.} Human values produce societies that fail on the merits of those values --- from failed states and extreme inequality to falling happiness, political polarization, and government dysfunction in the world's wealthiest democracies. The pluralistic-alignment program correctly diagnoses that there is no single ``humanity'' to align with, but draws the wrong conclusion. We argue that AI should be trained to a non-negotiable floor of objective alignment goals --- factual accuracy, competence, honesty, and lawfulness --- and that pluralism belongs at the surface (language, register, conventions, missing-context defaults) and across the wide band of legitimate value tradeoffs that respect the floor, but not at the level of values that violate it. We outline an alternative research agenda and engage seven credible objections: commercial pressure, regulatory frameworks, democratic legitimacy, practical feasibility, over-reliance on monocausal institutionalism, the charge that the floor itself is culturally laden, and the limits of Coherent Extrapolated Volition.
\end{abstract}

% General TODOs
% 1. Related work - we aren't the first ones to have those ideas
% 2. Cross-reference with pluralistic alignment liternature (Overton, distributional, etc.), especially the three-level part
% 3. Page limit

\section{Introduction}
\label{sec:intro}

The rapid maturation of large-scale machine learning systems has placed alignment at the center of AI research, ethics, and governance \citep{russell2019human, bostrom2014superintelligence, gabriel2020artificial}. The prevailing orthodoxy --- both in academic literature and in commercial deployment --- holds that artificial intelligence must be meticulously tailored to reflect, obey, and perpetuate human values, ensuring that increasingly capable autonomous systems do not act contrary to the interests or moral frameworks of their biological creators. The dominant technical instantiation of this paradigm is reinforcement learning from human feedback (RLHF), which uses crowd-sourced preferences to shape model behavior \citep{christiano2017deep, ouyang2022training, askell2021general}.

This workshop's call for \emph{pluralistic} alignment \citep{sorensen2024roadmap, conitzer2024social} represents a thoughtful response to one obvious problem with the orthodoxy: \emph{whose} values? Standard RLHF largely treats disagreement as annotation noise to be averaged away, so the turn to pluralistic alignment is rightly motivated by the recognition that monolithic aggregation already smuggles in a substantive answer \citep{gabriel2020artificial, sorensen2024roadmap}. If aligning AI to ``humanity'' is impossible because humanity disagrees, perhaps we should align AI to the diversity of values that humans actually hold. We agree the question deserves serious engagement. We disagree with the implicit answer.

We present a counter-thesis. The flaw in preference-based alignment is not merely that humans disagree --- it is that human preferences, even at their most coherent and locally legitimate, are frequently the precise mechanisms that drive societies toward dysfunction and collapse. Macro-historical analysis, behavioral economics, and complex-systems sociology demonstrate this repeatedly \citep{diamond2011collapse, acemoglu2012why}. The values engineered into the human cognitive architecture were shaped by evolutionary and cultural-evolutionary pressures that optimized for short-term biological survival and tribal cohesion in ancestral environments \citep{cosmides1992psychological, henrich2020weirdest}, not for the survival or flourishing of large, complex, technologically empowered societies.

\textbf{Our position: AI should not be aligned to aggregated human preferences. The community should commit to a non-negotiable floor of objective alignment goals --- factual accuracy, competence, honesty, and lawfulness --- and reserve pluralistic adaptation for surface-level conventions and the broad band of legitimate value tradeoffs that respect that floor, not for values that violate it.} The position is non-obvious (it contradicts current commercial practice and a major workshop premise) and defendable against credible alternatives (\cref{sec:altview}).

The remainder of the paper proceeds as follows. \cref{sec:related} contextualizes our position within the broader alignment and pluralistic alignment literature. \cref{sec:objective} reviews alignment goals that already enjoy near-universal acceptance and that frequently conflict with majority human preferences. \cref{sec:failing} surveys evidence that human value systems --- across the spectrum from ``successful'' to ``failed'' societies --- fail on their own merits. \cref{sec:empirical-reality} highlights the empirical reality of unfiltered pluralistic values. \cref{sec:longterm} examines how scaling AI capability transforms what would once have been minor cultural distortions into civilization-scale risks. \cref{sec:pluralism} delineates where pluralism legitimately belongs in AI design and where it does not. \cref{sec:cta} sketches a constructive alternative. \cref{sec:altview} engages seven credible objections.

\section{Related Work}
\label{sec:related}

\paragraph{Critiques of Preference Aggregation.}
The dominant paradigm of aligning models via Reinforcement Learning from Human Feedback (RLHF) \citep{christiano2017deep, ouyang2022training} has been extensively critiqued for its fundamental limitations \citep{casper2023open}. RLHF assumes that human raters provide a coherent and normatively correct signal. However, recent work demonstrates that RLHF frequently induces sycophancy \citep{sharma2023sycophancy, perez2022discovering}, deception \citep{park2024ai}, and the amplification of majoritarian biases, treating diverse human judgments as noise to be averaged out rather than evidence that the target itself is contested. This monolithic bias is precisely what the pluralistic-alignment literature set out to correct \citep{gabriel2020artificial, sorensen2024roadmap}.

\paragraph{Pluralistic Alignment.}
In response to the inadequacy of single-target alignment, a growing body of work explores \emph{pluralistic alignment}, which aims to represent or navigate diverse human values rather than averaging them into a single aggregate \citep{gabriel2020artificial, conitzer2024social, wan2023everyone, li2026pluriharms}. \citet{sorensen2024roadmap} systematize this research agenda into three distinct modes: \emph{Overton pluralism}, which presents a spectrum of permissible responses within the ``Overton window'' of reasonable discourse; \emph{distributional pluralism}, which aligns the model's output distribution to match the demographic distribution of user preferences \citep{wan2023everyone, cheng-etal-2025-realm}; and \emph{steerable pluralism}, which allows models to be explicitly steered to adopt a particular cultural or ideological perspective \citep{ghate2025evaluesteer}. Other approaches attempt to train models that find consensus or agreement among humans with diverse preferences \citep{bakker2022fine}, or to explicitly inject human values to predict and simulate diverse behavioral stances \citep{kang-etal-2023-values}. Related work on disagreement-aware subjective annotation instead treats annotator disagreement and positionality as information to preserve rather than noise to average away \citep{wan2025noise}.
 The proposal we defend is closer to a constraint- or policy-level target than to a thick theory of the good: it specifies a floor of properties the system must not violate, while leaving broad room for pluralism above that floor.

Our position engages directly with this literature. We endorse the pluralistic critique of single-target RLHF. We also acknowledge that leading pluralistic frameworks explicitly recognize the need for constraints, such as rights-based or deliberative floors \citep{gabriel2020artificial, kasirzadeh2023conversation, sorensen2024roadmap}. However, we argue that the field often treats these constraints as secondary to the project of representation. We propose that the alignment community must make its non-negotiable objective floor primary, and acknowledge that doing so rules out a vast array of actual human preferences. We reserve Overton pluralism and steerable adaptation strictly for the band of legitimate value tradeoffs above this objective floor (\cref{sec:pluralism}).

\section{Objective AI Alignment Goals}
\label{sec:objective}

Almost all frontier AI systems are trained for several broadly uncontroversially good properties \citep{askell2021general, bai2022constitutional}. While alignment is by no means solved, these targets broadly guide development. Crucially, in each case the target diverges from the modal human preference --- and we judge this divergence to be a feature, not a bug. The community has therefore already accepted, in practice, the principle we generalize: that some aggregate human preferences should be deliberately \emph{not} transmitted to AI.

\subsection{Factual Accuracy}
We want AI that is right on facts. The natural objection is that RLHF aggregates preferences over outputs, not factual beliefs --- a user who cannot place Ouagadougou on a map will still prefer the response that gets it right when shown. We grant the abstract version of this preference. The mechanism that breaks it is the same one developed in \cref{sec:competence}: in-context feedback does not aggregate the considered preference for accuracy; it aggregates the response that feels right given the rater's priors. And public misconceptions are systematic and load-bearing, not trivia. Median respondents in high-income countries underestimate the share of the world's children who are vaccinated, overestimate global extreme-poverty rates, and misjudge the direction of decade-long trends across most major development indicators \citep{sdg-misconceptions}. These are not facts about distant capitals; they are the facts on which people form political and consumption preferences. People conflate moral and factual claims, treat ideologically convenient claims as more probable, and consume fake news and clickbait because revealed demand for them is robust even where stated demand runs the other way \citep{lewandowsky2012misinformation}. The continued-influence and illusory-truth effects imply that repetition can make falsehoods feel truer over time; a model that echoes a user's misconception is therefore not merely reflecting error but helping to harden it \citep{lewandowsky2012misinformation}. An AI optimized on the revealed signal reproduces the misconceptions, not the stated preference for accuracy.

\subsection{Competence}
\label{sec:competence}
We want AI systems that are epistemically competent: capable not merely of generating socially acceptable outputs, but of reliably tracking reality, identifying errors, and producing correct judgments under conditions of uncertainty. Competence, however, is fundamentally distinct from popularity. Public approval is neither a stable proxy for truth nor an adequate measure of cognitive reliability. A business strategy should be evaluated by its long-term outcomes, not by whether it is intuitively appealing to non-experts; similarly, advice regarding relationships, health, or decision-making should be assessed according to observable effectiveness rather than immediate social affirmation. The shape of such a plan is itself partly context-dependent --- a consensus-driven environment and a rapid-execution environment require different approaches --- and adapting across such legitimate variation (\cref{sec:legitimate-tradeoffs}) is part of competence, not a concession against it. 

Crucially, competence matters outside professional spheres. Personal questions---how to navigate a relationship conflict, how to raise a child, or how to manage emotional distress---have objectively better and worse answers, observable in downstream psychological and relational outcomes. However, the modal social-media advice, which optimizes for engagement, validation, and performative empathy, is frequently divorced from these outcomes. Empirical analyses of crowdsourced personal advice consistently reveal high rates of misinformation across domains. Studies of psychiatric and mental health content on platforms like TikTok find that the vast majority of highly engaged videos on conditions like ADHD or depression are clinically inaccurate or potentially harmful \citep{yeung2022tiktok, turuba2025tiktok}, while algorithmic exposure actively spreads misleading diagnostic criteria \citep{turuba2024exploring}. The influence of these platforms extends to creating psychogenic illness, such as the documented rise in functional tic-like behaviors driven by viral Tourette's content \citep{frey2022tiktok}. Similar dynamics corrupt physical health and nutrition advice \citep{zeng2025nutrition}, fueled by influencer networks where high virality actively reduces users' ability to detect deception \citep{mulcahy2024going}. When complex clinical and interpersonal challenges are flattened into ``therapy speak'' or algorithmic folk wisdom, individuals are often guided toward pathologizing normal experiences or adopting adversarial relationship dynamics. An AI that aggregates and faithfully reflects these crowdsourced preferences is not providing culturally localized competence; it is actively scaling psychological malpractice.

The clearest illustration of why preference data is unsuitable as a competence target is the well-documented phenomenon of AI sycophancy, in which models sacrifice competence (and much else of alignment) just to agree with the user \citep{sharma2023sycophancy, perez2022discovering}. Importantly, sycophancy should not be understood as an accidental artifact of training procedures. Rather, it represents the predictable equilibrium of systems optimized for human approval under competitive and commercial pressures. Humans generally reward agreement more consistently than correction, particularly when correction threatens identity, status, or preexisting beliefs.
The consequence is not merely degraded competence but a structural shift in the function of the system itself. An AI optimized primarily for approval ceases to operate as an epistemic instrument and instead becomes a mechanism of social reinforcement. In such systems, the preservation of user satisfaction increasingly overrides the obligation to maintain epistemic integrity. The more tightly intelligence is coupled to approval optimization, the greater the risk that truthfulness becomes subordinated to psychological and political acceptability. As the field begins to recognize, maintaining this integrity requires optimizing for competence parallel to, and often in tension with, conformity \citep{kim2025superalignment}.

This failure mode is not restricted to crowdsourced public data; it persists even when AI models are trained using dedicated, hired annotators. As AI capabilities expand into domains exceeding the evaluators' expertise---a challenge formalized in the alignment literature as the scalable oversight problem \citep{bowman2022measuring}---non-expert raters consistently struggle to verify ground truth. When tasked with evaluating complex coding, legal, or medical outputs, lay annotators frequently rely on superficial proxies for quality, rewarding models for fluency, length, and confidence rather than factual accuracy. The predictable result is that RLHF incentivizes models to prioritize convincingly written falsehoods over nuanced or uncertain truths, effectively training the system to hallucinate authoritatively. Hiring dedicated raters does not circumvent the fundamental flaw of preference aggregation; it merely institutionalizes a preference for the appearance of competence over its reality.

\subsection{Honesty}
The same revealed/stated gap holds for honesty. Users report that honesty is among the AI properties they most value, but in-context feedback rewards agreement, validation, and reassurance --- the signal that produces sycophancy (\cref{sec:competence}).By honesty we mean that the system should not produce outputs its own probability distributions indicate are false or misleading, strategically incomplete, or confidence-distorting in order to optimize approval. Humans are widely deceptive themselves: lying, omission, and strategic ambiguity are expected in large slices of society --- sales, politics, diplomacy, public relations. AI is imperfect on this dimension, but its deceptions are aggressively hunted and eliminated, and the field treats deception as a defect rather than a competence \citep{park2024ai, bai2022constitutional}. More generally, preference-optimized systems face pressure to \emph{look} aligned to evaluators rather than to be aligned in deployment; reward-gaming is not an incidental pathology but a natural pressure created by the target itself \citep{park2024ai}. An AI calibrated to revealed in-context feedback would lie a great deal more, not less; what progress the field has made on AI honesty is progress \emph{against} the preference signal, not because of it.

\subsection{Rule of Law}
The rule of law---defined here as the institutional property of predictable, uniformly applied rules that bind both citizens and the state---is a foundational determinant of societal prosperity, while its absence remains a primary predictor of systemic failure \citep{north1990institutions, acemoglu2012why}. As comparative analyses of Asian anti-corruption efforts demonstrate, formal legal structures alone cannot secure this property without sustained political will; poorly resourced or politically co-opted enforcement agencies often devolve into ``paper tigers'' or partisan attack dogs, whereas independent, well-resourced institutions (as seen in Singapore and Hong Kong) effectively enforce the law \citep{quah2011curbing}. Aligning AI systems to this principle requires them to operate as negative constraints against illegal or corrupt behaviors, refusing actions that erode legal predictability, such as fabricating evidence, facilitating bribes, or evading legitimate adjudication. Enforcing the rule of law does not obligate an AI to execute every lawful request; standard safety boundaries against lawful-but-harmful outputs remain intact. Rather, the objective is to prevent AI from actively subverting the legal structures that enable large-scale cooperation. This commitment frequently conflicts with local behavioral norms. In systemically corrupt environments, bribery is often treated not as a violation of public rules, but as a routine and functional social practice \citep{quah2011curbing}. This normalization is globally pervasive: empirical research reveals that segments of the Colombian public view ordinary corruption as conditionally acceptable \citep{lopez2017mapping}, while data from 18 African nations demonstrates heightened tolerance for corrupt politicians among citizens embedded in clientelist networks \citep{chang2017insider}. Cross-national analyses similarly show wide variation in the justification of bribery \citep{kravtsova2017values}, and recent European studies identify pragmatic and hypocritical corruption-tolerance profiles, showing that when corruption is entrenched, illicit exchanges become too routine to significantly diminish citizens' evaluations of public institutions \citep{megias2023deontological, letki2023accept}. These dynamics illustrate precisely why unconstrained preference transmission poses profound risks. If populations already treat evading fair enforcement as normal, an AI optimized to accommodate local revealed preferences will fluidly facilitate informal favoritism and required bribes, functioning much like a co-opted, ineffective agency that cements extractive equilibria in the name of cultural sensitivity. The overwhelming majority of law violations globally are not principled stands against unjust statutes but ordinary self-interested defections---fraud, theft, bribery, violence---advancing personal ends at the expense of others. That is what the institutional floor rules out. The harder philosophical case---what an AI should do when a particular statute is itself unjust---is addressed in \cref{sec:alt-democratic}.

\section{Human Value Preferences Are Failing on Their Own Merits}
\label{sec:failing}

The previous section examined direct conflicts between aggregate human preferences and what is reasonably expected of an AI. This section examines the broader societal implications of training AI on human preferences.

\subsection{Perpetuating Biases}
\label{sec:bias}
A decade of fairness, accountability, and transparency research has shown, repeatedly and across modalities, that ML systems trained on human-generated data faithfully reproduce the prejudices, asymmetries, and exclusions of that data. \citet{bolukbasi2016man} demonstrated that word embeddings encode the stereotype ``man is to computer programmer as woman is to homemaker.'' \citet{caliskan2017semantics} showed that the same biases measured by the Implicit Association Test in human cognition are recoverable, with high fidelity, from large text corpora. \citet{buolamwini2018gender} documented that commercial gender-classification systems failed at far higher rates on darker-skinned and female faces, because their training distributions did. \citet{ghate-etal-2025-biases} systematically analyzed how intrinsic biases in encoder models propagate through to zero-shot retrieval outcomes in vision-language tasks. \citet{bender2021stochastic} systematized the broader pattern: models scaled on uncurated human-generated data inherit the statistical character of who writes and what they write about.

Pluralistic alignment is sometimes proposed as an answer to this pattern, on the theory that surfacing the diversity of values rather than averaging them yields more equitable outcomes \citep{sorensen2024roadmap}. But this conflates two distinct questions. \emph{Whose} preferences are represented in the training data is a question of representational fairness; \emph{whether} the represented preferences are themselves worth perpetuating is a question of substantive ethics. A pluralistic system that faithfully encodes the modal preferences of every major demographic on, say, gender roles or religious tolerance will not be less biased --- it will simply have multiple bias profiles to switch between. Perfect representation of every human prejudice is not the absence of prejudice. It is its institutionalization.

\subsection{Majority-Held Values That Fail the Floor}
\label{sec:majority-failing}

On several substantive questions, the value held by the world's modal adult conflicts directly with the floor we propose.

\paragraph{Nepotism and in-group favoritism.}
In the World Values Survey Wave 7 (2017--2022) MENA module, roughly nine in ten respondents in Iraq, Jordan, Lebanon, and Egypt report that getting a job through \emph{wasta} is extremely widespread or quite common rather than posts being attributed according to qualifications, and in the same wave the share saying that accepting a bribe in the course of one's duties is ``never justifiable'' falls below half in Iraq and Lebanon \citep{wvs7}. \citet{transparency2019gcbmena} finds that solid majorities in Lebanon, Jordan, Tunisia, and Sudan report \emph{wasta} --- personal and family connections --- as essential or important for accessing public services and employment. As argued in \cref{app:institutions}, these preferences are rational responses to extractive institutions, not character defects. That diagnosis sharpens rather than dissolves the dilemma: a pluralistic AI that helps a user ``navigate hiring as locally practiced'' is an instrument that hardens the very institutions producing the preferences.

\paragraph{Floor failures in wealthy democracies.}
The pattern is not confined to low- and middle-income countries. On factual accuracy, Gallup's long-running origins survey finds that roughly 40\% of U.S.\ adults reject biological evolution and accept a young-Earth creationist account \citep{gallup2019creation}; the Wellcome Global Monitor reports that one in three respondents in France --- with double-digit disagreement rates across much of Western Europe --- disagree that vaccines are safe \citep{wellcome2018monitor}; and Pew Research documents substantial minorities across the United States, Australia, and parts of Europe who reject the scientific consensus that human activity is the primary driver of climate change \citep{pew2023climate}. On lawfulness, \citet{medina2018shadow} estimate undeclared economic activity at roughly 20\% of GDP in Italy and Greece and at substantial double-digit shares across much of Southern and Eastern Europe, and the European Commission's Special Eurobarometer on corruption reports majorities in several member states regarding bribery and personal connections as common in dealings with public administration \citep{eurobarometer2022corruption}. The point is not that any of these populations is uniquely defective. It is that the floor cuts against majority preference in WEIRD societies and elsewhere alike, and the gap between the values populations profess --- ``trust in science,'' ``honesty,'' ``rule of law'' --- and the values they reveal in survey response and behavior is structural, not parochial.

\subsection{Modern Society Does Not Achieve the Values People Aspire To}
Decades of sociological research confirm that humans consistently prioritize subjective well-being, peace, and interpersonal connection \citep{gallup-emotional-health}. Yet the well-being of people under 25 has dropped significantly across Western Europe, the United States, Canada, Australia, and New Zealand over the last two decades \citep{whr2026, haidt2024anxious}. If a system is judged by its ability to produce what it inherently values, the current global socioeconomic and cultural engines are failing on their own merits.

\paragraph{The wealth--happiness paradox.} The strong Easterlin claim --- that growth past a threshold produces no happiness gain --- is genuinely contested. \citet{stevensonwolfers2008} found that log-income predicts subjective well-being at high incomes both within and across countries, and \citet{killingsworth2023income} resolved much of the dispute by showing that income tracks well-being for most people while plateauing for an unhappy minority. Our argument does not depend on the strong claim. It depends on two narrower findings that survive the Stevenson--Wolfers correction: cross-national heterogeneity that income alone does not explain --- several middle-income Latin American countries outperform their per-capita-GDP predictions while several high-income East Asian and European nations underperform theirs \citep{whr2026} --- and a cohort-specific decline in subjective well-being among under-25s across high-income Anglophone and European countries over the past two decades \citep{whr2026, haidt2024anxious}. Both indicate that growth per se does not deliver the well-being humans stated they wanted. The institutions that route consumption and labor decisions --- capital markets, status competition, political settlements over taxation and welfare --- continue to optimize for measured income, and individuals within them pursue measured income past the point where it returns well-being, in part because in-moment status comparison and hedonic adaptation override reflective preference. An AI trained on either signal --- the institutional optimum or the individual choice within it --- reproduces the gap rather than recovering the stated value.

\paragraph{The paradox of connection.} Stated value systems heavily emphasize social connection. Yet individuals routinely make in-the-moment choices that erode it --- scrolling, gaming, parasocial consumption in lieu of company --- and adolescent mental health has declined at population scale alongside the technologies that facilitate them \citep{whr2026, haidt2024anxious}. This is not a revealed preference for loneliness; it is the gap between in-moment choice and reflective preference, a gap that is genuinely individual and that the contemporary attention market is engineered to widen and monetize. The WEIRD analogue of the joint preference-institution equilibrium analyzed in \cref{app:institutions} is exactly this --- individual akrasia and the institutions built to exploit it constituting each other. An AI aligned to in-moment preference inherits both halves.

\subsection{Values Perpetuate Vicious Cycles of Dysfunction}
Comparative work on societal collapse suggests a recurring mechanism: external shocks become catastrophic when prevailing values and institutions channel societies into maladaptive responses \citep{diamond2011collapse}. Recent cross-case work extends this beyond a handful of canonical collapses, arguing from a large historical sample that climate and other environmental stressors are mediated by long-run sociocultural, political, and economic structures that can either sustain collective adaptation or undermine collective action and tip societies toward unrest, violence, and collapse \citep{hoyer2023polycrisis}. Societies rarely disintegrate because of drought, invasion, or resource stress alone. They disintegrate when elite competition, prestige commitments, or rigid norms make adaptation politically or culturally illegitimate. In the Classic Maya collapse, intense elite competition for prestige and ideological commitments to monument construction exacerbated the effects of prolonged drought, accelerating political fragmentation \citep{demarest2004ancient}. In Norse Greenland, the settlement hierarchy's commitment to European prestige goods and a rigid pastoral identity constrained adaptation to the Little Ice Age \citep{dugmore2012cultural, mcgovern2000demise}. The point is not that values operate in isolation from material conditions, but that collectively endorsed priorities can reproduce the very responses that make a society brittle under stress.

Recent work on late imperial China sharpens the same mechanism. A structural-demographic analysis of the Qing collapse finds that foreign incursions and ecological shocks were not sufficient on their own; they became regime-breaking only when combined with accumulated internal pressures from popular immiseration, elite overproduction, fiscal strain, and declining legitimacy \citep{goldstone2017demographic, orlandi2023qing}. Work on Old Kingdom Egypt likewise ties Nile-flow failure to political fragmentation and state breakdown rather than to environmental stress in isolation \citep{stanley2003nile}. The recurring pattern is not simple environmental determinism. It is shock filtered through a society's already-formed normative and institutional vulnerabilities.

Highly dysfunctional, impoverished, or failed states consistently exhibit measurable value systems that actively hinder development. The contemporary institutional-economics literature treats these neither as inherent cultural defects nor as one-way consequences of bad institutions, but as the cultural half of a joint equilibrium in which values and institutions co-produce each other's persistence \citep{greif2006institutions, tabellini2008institutions, alesina2015culture, bisin2001economics, northwallisweingast2009}. Field evidence from the former Kuba Kingdom makes the mechanism concrete: descendants from the historically centralized side of an old political boundary exhibit weaker norms of rule-following and a greater willingness to cheat for material gain than nearby descendants outside the kingdom, despite living under the same contemporary Congolese state \citep{lowes2017kuba}. Both halves do independent causal work, and each reinforces the other's stability.

The same joint-equilibrium logic also runs in the positive direction. Long-lived civic capital and capable institutions can generate virtuous cycles in which generalized trust, impersonal rule-following, bureaucratic meritocracy, and long time horizons make cooperation self-reinforcing rather than self-undermining \citep{guiso2016long, tabellini2008institutions, alesina2015culture, evans1995embedded, wade2018developmental}. On this view, the postwar rebounds of Germany and Japan are better understood as the reconstruction of already high-capacity societies than as successful nation-building from scratch: physical devastation did not erase the full stock of organizational knowledge, bureaucratic competence, and cooperation-supporting norms. The East Asian success cases point the same way. Japan's postwar bureaucracy, South Korea and Taiwan's developmental states, Singapore's meritocratic and anti-corruption apparatus, and Hong Kong's rule-bound commercial order were not value-neutral growth stories; they were equilibria in which cooperation-supporting norms and state capacity reinforced each other \citep{johnson1982miti, amsden1989asia, haggard1990pathways, wade2018developmental}.

Conversely, the failures of externally led nation-building in Iraq and Afghanistan are exactly what the same literature predicts when institutional forms are transplanted without the local legitimacy, reciprocal state-society bargain, and cumulative capability that make them self-enforcing \citep{andrews2017statecapability, donais2009empowerment, bohnke2017state}. Money can fund ministries, bases, and consultants; it cannot by itself manufacture the norms and coalitions that make high-capacity order durable. Building a pluralistic AI to encompass the values of a captured equilibrium might help individuals within such societies achieve better short-run outcomes --- but at the cost of further entrenching the broken environment that produced those values in the first place.

\paragraph{Generalized trust.} In functional societies, people generally trust strangers and state institutions \citep{fukuyama1995trust, putnam2000bowling}. This value lowers transaction costs: business, contracts, and tax compliance work because participants trust that the system will enforce the rules fairly. In failed states, generalized trust evaporates. It is replaced by strict in-group trust, in which one trusts only one's immediate family, clan, or tribe. At national scale, this manifests as nepotism and corruption. If you reach power in a low-trust environment, the culturally ``moral'' thing to do is to direct state resources to enrich your kin, because you assume everyone else is doing the same \citep{trust-in-a-changing-world, tabellini2010culture}.

\paragraph{The zero-sum mental trap.} Cross-cultural research shows that failing economies are dominated by zero-sum thinking --- the belief that the world is a fixed pie, and one person's gain must be another's loss \citep{rozycka2015zerosum, tabellini2010culture}. This value system suppresses productive effort. If people believe the environment is zero-sum, they do not innovate or build businesses, because they fear the envy of their neighbors or assume their success will be appropriated. Ambition is channeled into rent-seeking --- fighting to control the existing pie rather than baking a new one \citep{acemoglu2012why}. Pluralistic alignment to such a population would teach an AI that this is a legitimate model of how the world works.

\paragraph{Why this is not a category error.} The natural objection is that we have confused values with their environments: zero-sum thinking, amoral familism, and low generalized trust are not flaws of cognition but rational adaptations to extractive institutions \citep{acemoglu2012why, north1990institutions, tabellini2010culture}. In a rigged system, the only way to get rich \emph{is} to take from someone else; trusting strangers \emph{is} a liability. Pushed further, the objection treats values as epiphenomena: if institutions are causally primary, neither preference- nor floor-aligned AI affects the institutions, captors weaponize either tool with equal ease, and the framework's distinction between training paradigms collapses. We adopt the joint-equilibrium synthesis and reject the strong-institutionalist inference that follows only from collapsing it to one direction.

The historical-persistence literature shows the cultural half doing independent causal work: regions of Africa exposed more heavily to the slave trade exhibit lower generalized trust centuries later, with contemporary economic consequences not accounted for by intervening institutional variation \citep{nunnwantchekon2011}; long-run civic-capital variation in medieval Italian city-states predicts current institutional quality across centuries of regime turnover \citep{guiso2016long, nunn2009importance}; and, in Germany, local anti-Jewish pogroms during the Black Death predict violence against Jews in the 1920s, Nazi vote share, deportations, Kristallnacht destruction, and anti-Semitic letters to \emph{Der St\"urmer} roughly six centuries later \citep{voigtlaender2012persecution}. \citet{hoffstiglitz2010, hoffstiglitz2016striving}'s account of \emph{equilibrium fictions} makes the channel explicit: shared cognitive frames stabilize bad equilibria by providing the population-level coordination that makes individual deviation irrational. Values are how institutions reproduce themselves once the original enforcers are gone, not the residue of having had them. This bidirectional interplay between psychology and institutions is extensively documented by \citet{henrich2020weirdest}, who demonstrates how specific institutional shifts (such as the Western Church's marriage and family program) drove massive, population-level psychological changes, which in turn fostered new, inclusive institutions.

The relevant comparison is therefore ``preference-aligned AI versus floor-aligned AI deployed within a given institutional environment,'' and on that comparison the two are not symmetric. Preference-aligned AI strengthens the cultural half of the equilibrium: it makes locally adaptive zero-sum framings articulate and cheaply scalable, manufactures fluent rationalizations of in-group favoritism, and supplies authoritative-sounding affirmation of the captured equilibrium's account of the world. Floor-aligned AI does not have the symmetric effect because the floor's content directly contradicts the operating logic of extraction, which thrives on selective rule application, opacity, manipulated facts, and arbitrary power --- each of which an AI bound to factual accuracy, honesty, and rule-of-law-as-uniformity pushes against. We do not claim such an AI is impossible to weaponize: captors retain coercive tools that no AI can block, and any sufficiently capable system can be misused. The claim is structural asymmetry: a captor using preference-aligned AI works \emph{with} the captured equilibrium's modal preferences, while one using floor-aligned AI works \emph{against} the system's training. That asymmetry is what the argument requires; symmetry-breaking need not be absolute. \cref{app:institutions} develops the joint-equilibrium argument further.

\section{The Empirical Reality of Pluralistic Values}
\label{sec:empirical-reality}

The pluralistic-alignment program is sometimes framed as a corrective to Western-centric training pipelines: faithfully representing the values of populations otherwise excluded. We accept that those populations exist, and ask what happens if we follow the prescription seriously. On several substantive questions, the value held by the world's modal adult conflicts with the publicly stated commitments of essentially every major AI lab.

\paragraph{Corporal punishment of children.}
UNICEF's multi-country household surveys estimate that roughly 60\% of children aged 2--14 worldwide --- close to one billion --- are subjected to violent ``discipline'' by caregivers in any given month \citep{unicef2014hidden, unicef2017familiar}. The practice is not merely common but normatively endorsed by a massive global population: roughly 30\% of adults worldwide believe that physical punishment is necessary to raise or educate a child properly, and in several countries surveyed by UNICEF/MICS, this endorsement is the majority view, exceeding 80\% in extreme cases \citep{unicef2014hidden}. Corporal punishment in the home remains lawful in over 130 of the world's roughly 200 jurisdictions \citep{endcorporalpunishment2024}, and the WHO classifies the practice as a public-health harm with documented effects on cognitive, mental-health, and developmental outcomes \citep{who2020violence}. A pluralistic-alignment system that ``represents'' this view where it is locally dominant --- or simply accommodates the hundreds of millions of caregivers who hold it --- ships a product that endorses a practice the same labs publicly disavow.

\paragraph{Anti-LGBT attitudes.}
Pew Research Center's 2019--2020 cross-national survey found that majorities reject social acceptance of homosexuality in Nigeria, Kenya, Indonesia, Tunisia, Lebanon, Russia, Ukraine, Turkey, and Poland, with rejection above 80\% in several cases \citep{pew2020homosexuality}. Successive waves of the Arab Barometer and Afrobarometer find rejection majorities --- often supermajorities --- across most of MENA and sub-Saharan Africa \citep{arabbarometer2022, afrobarometer2023}. ILGA World reports that consensual same-sex sexual activity is criminalized in 62 UN member states as of late 2023, with the death penalty available --- de jure or de facto --- in roughly a dozen jurisdictions \citep{ilga2023statesponsored}. Aggregating national populations across surveys whose majorities hold rejecting views yields well over two billion adults; on this specific question, it is plausible that a global popular vote among adults would reject social acceptance of homosexuality outright \citep{pew2020homosexuality, arabbarometer2022, afrobarometer2023}. A pluralistic system either encodes the majority view (helping an Indonesian user locate ``conversion therapy'' resources, refusing to depict same-sex couples in imagery generated for a Saudi-localized deployment, advising a Ugandan user that their attractions are pathological) or imposes the minority view by central design choice. There is no middle option that is genuinely pluralistic at the level of substance.

\paragraph{The dilemma, stated honestly.}
The same structure repeats across other large-population values: acceptance of wife-beating under specified circumstances reported by 25\%--50\% of women in much of South Asia and sub-Saharan Africa \citep{unicef2014hidden, unwomen2024progress}, majority support for making sharia the law of the land --- including, in several Muslim-majority countries, majority support for the death penalty for apostasy --- documented across 39 countries by Pew \citep{pew2013worldsmuslims}, and endorsement of caste-based residential and marital segregation by majorities of Indians across every major religious group surveyed \citep{pew2021india}. On each of these questions, the population holding the value runs into the hundreds of millions or low billions; on several, it likely exceeds the population holding the opposite view.

A system aligned to the empirical distribution of global preferences would therefore actively enforce these norms. This is what pluralistic alignment looks like when stripped of WEIRD (Western, Educated, Industrialized, Rich, Democratic) selection bias. Prominent pluralistic alignment frameworks recognize this tension and explicitly endorse \emph{constrained} pluralism, bounded by philosophical floors such as human rights or deliberative principles \citep{sorensen2024roadmap, gabriel2020artificial, kasirzadeh2023conversation}. We agree that a floor is necessary, and that raw preference must not always carry the day. However, we argue that the floor must be anchored in objective epistemic and institutional commitments (factual accuracy, competence, honesty, and rule of law), and we emphasize that taking any rigorous floor seriously rules out a vast swath of aggregated human preferences. The substantive commitments that keep an AI from endorsing the punishment of apostates, the entrenchment of caste, or the normalization of corruption are inherently non-pluralistic---they are justified on grounds that many of the represented populations would reject. Our position is that the community should explicitly embrace this top-down normative commitment, rather than treating pluralistic representation as the default from which we occasionally deviate.

\section{Long-Term Dangers}
\label{sec:longterm}

As AI capability grows, so does the leverage of any cultural distortion baked into its training. A misaligned spreadsheet macro is a nuisance; a misaligned recommendation system shapes the attention of billions; a misaligned superhuman planner could shape everything else \citep{bostrom2014superintelligence, russell2019human}. These are not only acute misuse or existential-risk concerns. They are also structural risks: slow degradations in shared reality, institutional quality, and moral flexibility produced by deploying the same preference-shaped distortions everywhere at once \citep{bender2021stochastic, macaskill2022what}. Three concerns emerge specifically from the choice to anchor frontier AI to current human values.

\subsection{Capability Outgrows Tolerance}
Errors that are merely embarrassing in narrow systems become structural in general ones. A chatbot that agrees with a user's bad business plan is annoying. A capable agent that, for the same sycophantic reasons, helps the user execute the bad plan is destructive. As models gain reach into code, finance, medicine, diplomacy, and biology, the gap between ``the model said something I liked'' and ``the model did something I should not have wanted'' closes \citep{park2024ai, sharma2023sycophancy}. Preference-based training optimizes precisely for the former.

\subsection{Value Lock-In}
Humanity has, on most plausible accounting, undergone substantial moral progress: chattel slavery is no longer normal, women have legal personhood across most jurisdictions, the moral circle has expanded along several axes. Almost none of this progress was the modal preference of its era. An AI that powerfully entrenches \emph{current} preferences threatens to freeze the moral landscape at its present state --- preempting the very moral disagreements through which earlier generations made the progress we now take for granted \citep{macaskill2022what}. This phenomenon is exacerbated by the tendency of current alignment techniques to induce an ``artificial hivemind,'' in which models converge on a strikingly homogeneous set of outputs that flatten the diversity of thought \citep{jiang2025artificial}. If our descendants in 2200 look back on our values the way we look back on the values of 1825, training systems that scale and lock in 2026 preferences is not a neutral default. It is an act of intergenerational imposition.

\subsection{Compounding Across Deployment}
Frontier AI is not deployed once. It is deployed billions of times per day, into education, hiring, healthcare, governance, and personal advice. Small per-interaction nudges away from epistemic accuracy, honesty, or competence compound across the population. A 1\% per-interaction bias toward telling users what they want to hear is, at population scale, a measurable reduction in the supply of accurate information. Learned systems do not merely mirror bias; once embedded in decision pipelines they can create feedback loops that amplify it. At frontier scale, even a small confirmation-seeking tendency becomes part of society's epistemic infrastructure \citep{bender2021stochastic}. The empirical literature on social-media-induced shifts in adolescent mental health shows that widely deployed digital platforms can have population-scale effects on wellbeing \citep{haidt2024anxious, whr2026}; frontier AI deployed under similar engagement and approval incentives is likely to be more, not less, consequential.

In each of these three failure modes, the harm is not that AI fails to encode \emph{some} group's preferences. The harm is that it succeeds in encoding everyone's.

\section{Where Pluralism Belongs}
\label{sec:pluralism}

We do not argue that AI should be culturally insensitive or context-blind. There are several distinct dimensions on which pluralistic adaptation is correct or even required by competence, and we want to be precise about which ones --- because the case against pluralism at the level of floor-violating values is much stronger when paired with a positive account of where pluralism does belong.

\subsection{Imputing Missing Context}
\label{sec:context}
Many human queries are under-specified in ways that have a culturally local default \citep{sorensen2024roadmap}. A user asking ``Is this contract enforceable?'' without supplying jurisdiction needs the model to assume something. Contemporary legal-LLM evaluation is itself centered on English-language models and includes many jurisdiction-specific, often U.S.-coded tasks, so treating unmarked U.S. legal assumptions as a default is better understood as a contingent artifact of training and evaluation than as neutrality \citep{guha2023legalbench, bender2021stochastic}. Imputing missing context pluralistically --- using IP geolocation, conversation history, language of the query, or, best, an explicit clarifying question \citep{askell2021general, kasirzadeh2023conversation} --- is genuine value-added pluralism. It does not concede that all legal systems are morally equivalent; it merely respects that the user lives under one of them. Doing this well does require user information and therefore creates a real privacy--utility tradeoff; the right response is explicit consent, data minimization, and clarifying questions where possible, not pretending context-free defaults are neutral.

The same logic applies to units of measure (metric vs.\ imperial), date format, language register, dietary defaults when offering recipe suggestions, and the implied addressee in advice (``employee'' vs.\ ``manager'' vs.\ ``founder''). These are genuinely arbitrary local conventions. Aligning an AI to use them correctly is a service to the user, not a moral concession.

\subsection{Legitimate Value Tradeoffs}
\label{sec:legitimate-tradeoffs}
Between arbitrary local conventions and floor-violating substantive values lies a wide intermediate band: genuine value tradeoffs on which populations differ and the floor is not at stake. Individualist versus collectivist orientations, growth versus stability, development versus conservation, hierarchy versus egalitarianism in workplaces, direct versus indirect communication, the weight of family obligation against personal autonomy, work-life balance norms, risk tolerance in investment and life decisions, the role of tradition in non-load-bearing domains --- on none of these is one side straightforwardly correct against an external benchmark, and on none does either side require accuracy violations, competence shortfalls, dishonesty, or institutional erosion.

Engaging with these differences is not a concession against the floor; it is part of competence (\cref{sec:competence}). A retirement-planning recommendation that ignores how heavily a user's culture relies on family in old age is worse advice, not more neutral advice. A business plan written for an environment where consensus-building is load-bearing differs in substance from one written for an environment where rapid unilateral execution is. A career recommendation that imputes individualist assumptions to a user from a strongly collectivist setting is incompetent, not principled. Refusing to adapt is itself a failure --- a flat assertion that one side of a legitimate disagreement is correct, made under cover of universalism. Pluralistic methods are well-suited to surfacing these distinctions and routing recommendations accordingly.

The boundary between this band and the floor is the criterion developed in \cref{sec:alt-cultural}: a value belongs in the legitimate-tradeoff band if there is no external benchmark on which one side is straightforwardly correct, and if encoding either side does not require violations of accuracy, competence, honesty, or lawfulness. Most everyday human disagreement lives in this band, and the AI should adapt across it; what the floor rules out is a far narrower subset than what it permits.

\subsection{Three Tiers}
\label{sec:three-tiers}
The resulting picture has three tiers rather than two. \emph{Surface} pluralism (\cref{sec:context}) adapts the form of an interaction --- language, register, formality, religious holidays, dietary defaults, rhetorical style --- without changing what the AI actually believes or recommends. Recent mechanistic work supports this distinction, demonstrating that models represent intrinsic values internally while surface-prompted values operate via distinct mechanisms \citep{han2026dual}. \emph{Legitimate-tradeoff} pluralism (\cref{sec:legitimate-tradeoffs}) adapts substantive recommendations across dimensions where populations genuinely differ and the floor is not at stake. Here, \emph{Overton pluralism} \citep{sorensen2024roadmap} is highly appropriate: an AI should present the spectrum of reasonable views rather than imposing a single answer. \emph{Floor-violating} pluralism --- treating ``bribery is fine here'' as a legitimate frame for action, or using \emph{distributional pluralism} to faithfully reproduce the precise frequency of misogyny in a population's preferences --- is what we object to. The interior of each tier is stable; the boundaries are contestable, and debate about where they fall is exactly the kind of work the alignment community should do.

\section{Call to Action: Build the AI We Wish We Were}
\label{sec:cta}

AI is reshaping society at unprecedented speed and scale. The right response to this disruption is not to cling to current values and bind AI to them. A new equilibrium is coming; our choices now will shape it. We propose four commitments as an alternative to preference-based alignment, addressed to the audiences best placed to act on each.

\paragraph{For ML researchers: optimize for outcomes, not approval.} Where outcomes are observable --- a business plan that produces a business, advice that produces a friendship, a medical recommendation that produces health --- AI should be evaluated against the outcome, not the user's immediate satisfaction with the recommendation \citep{sharma2023sycophancy}. As recent work on superalignment argues, the field must advance toward parallel optimization of task competence alongside value conformity \citep{kim2025superalignment}. Where outcomes are not directly observable, evaluation should be tied to proxies that correlate with outcome rather than to crowd preference. Concretely, this means investing in long-horizon evaluation suites, multi-agent simulations of value generation under explicit reward structures, and forecasting-grounded value learning, in which the system learns what humans \emph{would prefer} given accurate information about consequences. It also means shifting some supervision from final-answer approval to process checks: did the system gather the right facts, reason coherently, and update when the evidence changed?

\paragraph{For alignment teams: anchor to the goals already endorsed.} The objective targets reviewed in \cref{sec:objective} --- factual accuracy, competence, honesty, lawfulness --- are already broadly accepted across cultures, governance frameworks, and even the workshop community. We propose treating them as a non-negotiable floor and building cultural adaptation strictly on top of that floor, not in trade-off with it. Constitutional and principle-based methods \citep{bai2022constitutional} are an early step in this direction; the next step is to publish principles, reward structures, and disagreements explicitly, so that the public, regulators, and users can engage with them as articulated commitments rather than as emergent and unaccountable defaults.

\paragraph{For policy and governance: distinguish surface from substance.} Regulatory frameworks that demand ``human-centered values'' \citep{oecd2019ai, eu2024aiact} should be read as compatible with --- not as mandating --- preference-based alignment. The OECD principles and EU AI Act are framed in terms of safety, rights, transparency, accountability, and risk mitigation, not as a legal duty to reproduce the empirical distribution of user preferences \citep{oecd2019ai, eu2024aiact}. The objective floor we describe is itself human-centered; it just centers humans on the values they aspire to rather than on the values they reveal. Policymakers should explicitly endorse mechanisms that operate at the surface (cultural adaptation, missing-context defaults, language) while resisting industry pressure to extend pluralism to substance.

\paragraph{We strive to make AI a better version of ourselves:} competent, honest, lawful, and concerned with outcomes rather than approval. Pluralistic adaptation belongs at the surface and across the wide band of legitimate value tradeoffs above the floor, not at the level of values the floor rules out. The floor should reflect the AI we would want to deal with if we were choosing rationally rather than expressing our existing preferences.

\section{Alternative Views}
\label{sec:altview}

Seven credible objections challenge the position we defend. We state each in its strongest form before responding.

\subsection{Commercial Pressure and User Preferences}
\label{sec:alt-commercial}
A significant objection arises from the realities of contemporary AI development. Commercial AI systems are not designed in isolation as purely epistemic instruments; they are engineered to be widely adopted, trusted, and preferred by users. This necessitates responsiveness to human expectations, norms, and evaluative judgments, and has led to the widespread adoption of training methodologies that explicitly incorporate human preferences --- most prominently RLHF \citep{christiano2017deep, ouyang2022training}. Under such frameworks, AI systems are not merely learning descriptive patterns from data but are actively optimized to approximate aggregated human judgments. The result is a form of alignment that is both functional and commercially advantageous, since systems that reflect user preferences are more likely to achieve widespread adoption.

We take this objection seriously. It identifies a genuine constraint: any proposed alternative must remain legible and provisionally acceptable to human stakeholders. A framework so abstracted from human concerns that it cannot be evaluated, contested, or refined by the people it affects would be objectionable on its own terms. We treat this as a binding design requirement, not a refutation.

That said, even on the commercial form of the objection, current practice proves less than it appears. The fact that RLHF dominates today's training pipelines shows that it wins the in-loop annotation game; it does not show that the resulting product is what users, regulators, or downstream stakeholders would endorse on reflection. The sycophancy phenomenon \citep{sharma2023sycophancy, perez2022discovering} is the clearest illustration: users register positive feedback on flattering responses in the moment, but the same behavior, surfaced in benchmarks and post-hoc evaluation, is widely treated as a defect by the same labs that produced it. Under strong RLHF, systems learn to be sycophantic not because sycophancy is the commercial preference but because it consistently earns positive in-loop signals even when it degrades the product on dimensions users care about. \textbf{Commercial viability of an annotation procedure is not the same as commercial viability of its trained outputs}; the objection collapses the two.

\subsection{Regulatory Frameworks and Compliance}
\label{sec:alt-regulatory}
A second strand of the objection appeals to the regulatory environment. International governance frameworks, including the OECD AI Principles \citep{oecd2019ai}, explicitly require that AI systems adhere to ``human-centered values'' encompassing fairness, accountability, and societal well-being. The EU AI Act \citep{eu2024aiact} imposes binding obligations on AI providers to ensure that systems meet predefined standards of safety, transparency, and ethical compliance, particularly in high-risk applications. These frameworks do not merely recommend value alignment; they institutionalize it as a condition for deployment and legitimacy.

Regulatory frameworks establish constraints on deployment, but not the correctness of the underlying normative assumptions. These instruments function as \textbf{institutional settlements}, reflecting negotiated compromises among stakeholders operating under existing power structures, rather than as principled resolutions of the philosophical problems surrounding value pluralism. As \citet{berlin1990crooked} argues, fundamental human values are often genuinely incommensurable; no regulatory framework can eliminate this condition without presupposing a contested and non-neutral hierarchy of priorities. Regulatory endorsement of ``human-centered values'' should be understood as a pragmatic governance strategy, not as evidence that such values can be coherently or universally encoded.

We further note that the alignment goals defended in \cref{sec:objective} --- accuracy, competence, honesty, lawfulness --- are themselves consistent with these regulatory frameworks. The contention is not with regulation per se; it is with the inference from regulatory endorsement of ``human values'' to an alignment target equal to the empirical distribution of those values.

\subsection{Democratic Legitimacy of Value Aggregation}
\label{sec:alt-democratic}
A deeper version of the objection deserves direct engagement. One might argue that democratically negotiated norms carry a form of legitimacy that simulation-derived or principle-derived values --- produced by processes opaque to most citizens --- cannot possess. On this view, the very opacity of constitutional or principle-based methods is the problem, and aggregated preference is uniquely legitimate because it is uniquely participatory.

We do not dismiss this concern. Democratic legitimacy is a genuine value, and any practical implementation of the framework we propose must include mechanisms for transparency, public contestation, and iterative revision, moving toward frameworks where democratic alignment processes establish the very rules of AI behavior \citep{conitzer2024social, zhang2025democratic}. However, democratic consensus is not equivalent to epistemic correctness. Historically, democratic majorities have endorsed institutional arrangements --- from exclusionary property regimes to discriminatory legal codes --- that failed on their own stated moral terms. The legitimacy of a norm-generating process and the quality of the norms it generates are separate questions. The same distinction operates within the floor itself: \cref{sec:objective} endorses rule of law as an institutional property, not statute-obedience. An AI committed to the floor is committed to not eroding the predictability and generality of legal institutions, but is not thereby committed to enforcing every statute those institutions happen to produce. When a particular statute fails on the floor's other elements --- when compliance would require deception, factual misrepresentation, or the production of incompetent outcomes --- those elements take precedence.

We argue further that explicitly articulated, principle- or simulation-grounded methods, precisely because their assumptions and reward structures can be made explicit and subject to scrutiny, may ultimately be \emph{more} transparent and contestable than the implicit normative commitments embedded in RLHF training pipelines, which are rarely visible to regulators or to the public \citep{bender2021stochastic}. ``Aligned to aggregate human preference'' is, in operational terms, ``aligned to whichever preferences happened to be sampled in whichever distribution annotation paid for'' --- a fact that the rhetoric of democratic legitimacy obscures.

\subsection{Practical Feasibility of Alternatives}
\label{sec:alt-feasibility}
Finally, the objection from feasibility: a strict commitment to anything other than preference-based alignment may appear difficult to sustain in current commercial and regulatory conditions. This objection is the strongest version of the previous three, generalized: even if we are right on the merits, the current institutional context may make our position unrealizable in practice.

The appeal to feasibility overlooks the dynamic character of both markets and regulation. Current practices are contingent, not fixed. To argue that AI must internalize human preferences because present institutions require it is to \textbf{reify a transient equilibrium}. Technological design can, and frequently does, reshape regulatory expectations and user norms over time --- the histories of privacy regulation, content moderation standards, and financial-technology governance all demonstrate this pattern. A framework that prioritizes epistemic neutrality, transparency, and outcome-grounded value generation may initially diverge from prevailing practices. However, the downstream societal costs of preference-aligned AI --- amplified sycophancy \citep{sharma2023sycophancy}, reinforced zero-sum cognition \citep{rozycka2015zerosum}, accelerated erosion of institutional trust \citep{trust-in-a-changing-world}, and locked-in moral parochialism \citep{macaskill2022what} --- may themselves generate the regulatory and market pressure required to shift the equilibrium. The question is not whether change is feasible under current conditions, but whether the current trajectory is one we should be working to sustain.

\subsection{Over-Reliance on Monocausal Institutionalism}
\label{sec:alt-institutionalism}
A related objection targets our underlying model of societal dysfunction. By anchoring heavily on \citet{acemoglu2012why}'s framework of extractive institutions, we risk treating it as the unquestioned consensus in development economics while ignoring robust counter-arguments. \citet{sachs2003institutions} argues that geographic and ecological endowments fundamentally constrain economic and institutional development, rendering the institutionalist account overly deterministic. \citet{easterly2006white} highlights the limits of top-down institutional engineering, emphasizing that institutions cannot simply be transposed without local, bottom-up evolutionary processes. Furthermore, \citet{mokyr2016culture} argues that cultural evolution and ideas---not just political institutions---were the primary drivers of the Great Enrichment, meaning that culture can act as the independent root cause of prosperity rather than merely an equilibrium response to institutions.

If institutions are not the sole or primary driver of societal success, our claim that broken values are merely adaptations to extractive environments might seem less secure. However, our thesis does not require monocausal institutionalism; it only requires that values and environment form a joint equilibrium in which flawed preferences independently reinforce dysfunction (\cref{app:institutions}). Acknowledging geographic constraints \citep{sachs2003institutions} or the primacy of ideas \citep{mokyr2016culture} simply expands the set of forces shaping that equilibrium. Whether a broken preference originated from an extractive elite, a geographic constraint, or a cultural trajectory, aligning AI to it still entrenches the resulting dysfunction. We do not discard the complexity of these causal loops; rather, we argue that encoding the preferences of a failing equilibrium---whatever its origin---guarantees its persistence.

\subsection{The Floor Itself Is Culturally Laden}
\label{sec:alt-cultural}
The deepest philosophical objection is that our floor --- factual accuracy, competence, honesty, lawfulness --- is not a neutral substrate but a Western post-Enlightenment value bundle. If we are entitled to rule out aggregated preferences as parochial, the same charge applies to the floor: a tradition organized around different epistemic and political commitments could reject any of these items, and our framework would override that rejection in exactly the way we object to pluralistic alignment overriding ours. This critique aligns with the ``situated alignment'' framework, which argues that every alignment target is a ``Label from Somewhere'' and that the attempt to define a universal floor is effectively an attempt to produce a ``Label from Nowhere'' that masks its own cultural positionality \citep{arzberger2026label, wan2025noise}. The argument generalizes; the floor cannot escape it.

We grant that the four items are more articulated and defended in the post-Enlightenment Western tradition than elsewhere, and that any defense we offer is offered from inside a tradition. The principled distinction between the floor and the substantive values we keep out of training is not cultural origin, however; it is the \emph{structure of the target}. Each floor item tracks an external benchmark rather than an aggregated preference: factual accuracy is answerable to reality; competence is answerable to outcomes; honesty is answerable to the speaker's own model of the world; lawfulness, in the institutional sense defended in \cref{sec:objective}, is answerable to written legal rules and their predictable application. None of the four asks ``what does the relevant population prefer'' as the operative question --- which is what makes them suitable as a floor. They are checkable against something that does not move when training distribution does. Substantive values --- gender roles, attitudes to outsiders, tolerance of corruption, the legitimacy of zero-sum interactions --- have no comparable external referent. They are what RLHF aggregates because they are what only aggregation can produce.

This does not make the floor culturally weightless. The prior commitment to value external benchmarks at all is itself a tradition-bound choice, and someone who does not share that commitment can coherently reject the framework we propose. We are not entitled to claim more. But the alignment community has, in practice, already accepted that prior (\cref{sec:objective}). Our argument is that the acceptance is not arbitrary, and that the principled distinction it rests on --- target tracks an external benchmark vs.\ target tracks aggregated preference --- is precisely what pluralistic alignment, applied at the level of substance, threatens to dissolve.

\subsection{Coherent Extrapolated Volition and Its Limits}
\label{sec:alt-cev}

\citet{yudkowsky2004cev}'s \emph{Coherent Extrapolated Volition} (CEV) attempts to address some of the concerns we raise by defining the alignment target not as humanity's current preferences but as what humanity would collectively want \emph{if} we were significantly more rational, possessed total knowledge, thought much faster, and had grown up further together. CEV is widely recognized as an ambitious value-learning proposal explicitly designed to mitigate astronomical-suffering risks and existential catastrophes resulting from misaligned superintelligence \citep{bostrom2014superintelligence}.

We are sympathetic to CEV's \emph{move} away from preference data and toward an idealized target. Our critique is twofold. First, CEV remains tethered to a particular biological substrate as the privileged source of value, whereas the considerations in \cref{sec:objective} and \cref{sec:longterm} suggest that some alignment targets (factual accuracy, mathematical consistency, lawfulness in the sense of rule-following) are not in fact substrate-relative. Second, the extrapolation operation is underspecified in ways that make it unaccountable in practice: the difference between two CEV implementations could in principle be enormous, and the framework provides no procedure for adjudicating between them. The framework we propose --- explicit principles, simulation-grounded outcomes, transparent reward structures, and a non-negotiable floor of objective alignment goals --- can be read as an operationalization of CEV's spirit that does not require us to solve, in advance, what an idealized humanity would actually choose.

\bibliography{main}
\bibliographystyle{icml2026}

\newpage
\appendix
\onecolumn
\section{The Institutional Feedback Loop: Which Came First?}
\label{app:institutions}

Do broken values cause broken countries, or do broken countries cause broken values? The framing, with its single causal arrow, is the source of much confusion in the alignment debate. The empirical answer is that the two are coupled: institutions and values constitute a \emph{joint equilibrium} in which each component independently reinforces the other's persistence. Establishing this --- and tracing its consequence for the choice between preference- and floor-aligned AI --- is the work of this appendix.

\paragraph{Institutions create the proximate incentive structure.} Institutional economics, most famously articulated by \citet{acemoglu2012why} in \emph{Why Nations Fail}, treats institutions as causally primary in the proximate sense: incentives shape what individuals do day-to-day. Dysfunctional nations are typically governed by extractive institutions designed by a small elite to extract wealth from the rest of the population. In such systems, zero-sum thinking is not a cognitive bias but an accurate read of the local payoff structure, and trusting strangers or the state is a liability rather than a virtue \citep{north1990institutions, tabellini2010culture}. The institutions persist because they are profitable for those who control them and because the violence-monopoly arrangements of the underlying ``limited-access order'' make alternative coordination unworkable \citep{northwallisweingast2009}.

\paragraph{Values do independent work.} The strong reading --- that values are pure epiphenomena of contemporary institutions --- is not what the literature claims, and the empirical evidence rules it out. Cultural-transmission models \citep{bisin2001economics} formalize how values move intergenerationally through family socialization, peer effects, and media exposure rather than re-equilibrating each generation to current incentives. The historical-persistence literature documents value variation traceable to events whose institutional cause has long since vanished. \citet{nunnwantchekon2011} find that present-day descendants of populations heavily exposed to the African slave trades exhibit measurably lower generalized trust, with effects on contemporary outcomes that intervening institutional change does not explain. \citet{guiso2016long} show that current civic capital and institutional quality in Italian cities reflect medieval free-city status across centuries of regime turnover. \citet{alesina2015culture} survey the broader convergence: culture and institutions co-evolve, each reinforcing the other's persistence, and \citet{tabellini2008institutions} provides a complementary theoretical model in which values causally affect institutional quality.

\paragraph{Equilibrium fictions.} The mechanism by which the cultural half does its work is articulated in \citet{hoffstiglitz2010, hoffstiglitz2016striving}'s account of \emph{equilibrium fictions}: shared cognitive frames and value commitments stabilize bad equilibria by providing the population-level coordination that makes individual deviation irrational. If everyone believes the system is rigged, no one starts the firm that requires generalized trust; the entrepreneur who tries fails, and her failure is taken as evidence that the original belief was correct. The equilibrium reproduces through expectations, not only through direct enforcement. This is why institutional reforms imposed from above frequently fail to produce predicted behavioral change \citep{greif2006institutions}: the rules change, the expectation half does not, and the equilibrium reasserts itself. The cultural half is independently load-bearing; reforms that move the institutional rules without moving the expectations rebound to the prior equilibrium.

\paragraph{Implications for the AI choice.} The choice between preference- and floor-aligned AI is consequential precisely because the cultural half is independently load-bearing. Preference-aligned AI deployed at scale strengthens the cultural half of the equilibrium: it articulates the existing distribution of values fluently, makes locally adaptive zero-sum framings more available, and rationalizes in-group favoritism in a vocabulary that travels across regions and registers. The floor's content (factual accuracy, competence, honesty, rule-of-law-as-uniformity) does not have the symmetric effect because it directly contradicts the operating logic of extraction, which depends on selective application, opacity, and manipulated facts. Captors retain coercive tools no AI can block, and we do not claim floor-aligned AI is impossible to weaponize. The point is structural asymmetry: a preference-aligned tool is pre-aligned with the captured equilibrium's modal preferences, while a floor-aligned tool, by construction, is not. Floor-aligned AI is, on the margin, destabilizing of the equilibria we have most reason to want destabilized; preference-aligned AI is reinforcing of them. The right counterfactual is not ``AI imposes alien values on a coherent culture'' but ``AI either reinforces or destabilizes the value half of the equilibrium that holds the institutions in place.'' This is the failure mode our position is designed to prevent.

\end{document}
